%------------------------------------------------------------------------------%
\begin{question}
  Analise o limite de \(f(x, y) = \frac{y}{x}\)
  quando \((x, y)\)  se aproxima de \((0, 0)\)

  \pause
  \vspace{\baselineskip}
  A reta \(x=0\) não faz parte do domínio de $f$

  \pause
  O ponto \((0, 0)\) não está no domínio de $f$

  \pause
  Não podemos usar as propriedades diretamente

  \pause
  Não existe manipulação que remova a indeterminação
\end{question}

%------------------------------------------------------------------------------%
\begin{resolution}{Solução}
  Analisando $f$ sobre a reta \(y=x\)
  \pause
  \[
    f(x, y) \evalat{y=x}{}
    \pause
    = f(x, x)
    \pause
    = \frac{x}{x}
    \pause
    = 1
    \pause
    \qquad\text{ pois } x=0 \text{ não está no domínio de } f
  \]

  \pause
  Analisando $f$ sobre a reta \(y=0\)
  \pause
  \[
    f(x, y) \evalat{y=0}{}
    \pause
    = f(x, 0)
    \pause
    = \frac{0}{x}
    \pause
    = 0
    \pause
    \qquad\text{ pois } x=0 \text{ não está no domínio de } f
  \]

  \pause
  O que vai acontecer com outas retas da forma \(y=ax\)?
\end{resolution}

%------------------------------------------------------------------------------%
\begin{resolution}{Solução}
  Para qualquer $\delta>0$
  \pause
  sempre vai existir um

  ponto da forma \((x, x)\) e um ponto da forma \((x, 0)\)

  \pause
  dentro da bola de raio $\delta$ de centro em \((0, 0)\)

  \pause
  Nenhum valor $L$ atende as condições da definição de limite

  \pause
  \vspace{\baselineskip}
  \emph{O limite não existe}
\end{resolution}

%------------------------------------------------------------------------------%
